% Môn: Toán
% Lớp: 10
% Chương: Chương I. Mệnh đề và tập hợp
% Bài:  Mệnh đề · Mệnh đề chứa biến
% Số câu: 0
% App đọc file này trực tiếp — sửa rồi Commit trên GitHub
% Ghi nguồn từng câu: \nguon{SGK} hoặc \nguon{2 SBT VL 12 KNTT} trong


%%%=============EX_1=============%%%

\dangbt{Mệnh đề chứa biến}
\begin{ex}
% Mức: TH
	Trong các câu sau, câu nào \textbf{không phải} là mệnh đề chứa biến?
	\choice
		{$x^2-2x-3=0$}
		{$3n+2\vdots 5$}
		{\True Số $1$ có một ước duy nhất là chính nó}
		{$x-5y=10$}
	\loigiai{
		Mệnh đề chứa biến là một câu khẳng định mà trong đó có chứa biến số, và tính đúng sai của nó phụ thuộc vào giá trị của biến số đó.
		Xét các phương án:
		Phương án A:  Đây là một mệnh đề chứa biến $x$. Tính đúng sai của nó phụ thuộc vào giá trị của $x$.
		Phương án B:  Đây là một mệnh đề chứa biến $n$. Tính đúng sai của nó phụ thuộc vào giá trị của $n$.
		Phương án C:  Đây là một câu khẳng định cụ thể, không chứa biến. Tính đúng sai của nó không phụ thuộc vào bất kỳ biến nào (câu này là một mệnh đề đúng).
		Phương án D:  Đây là một mệnh đề chứa hai biến $x, y$. Tính đúng sai của nó phụ thuộc vào giá trị của $x$ và $y$.
		Vậy, câu không phải là mệnh đề chứa biến là phương án C.
		}
\end{ex}